% answers_short.tex —— 2010 年·普通高等学校招生考试·安徽卷·理科 · 简短答案速查
% 由 tools/make_short_answers.py 自 answers.tex 的 \gksolhead/\gkanswer 提取，
% 并与 2010/2010_数学真题_安徽理_答案要点.md 逐题比对后生成，勿手改（改源 MD 后重跑）。

\begin{gknotice}
\textbf{使用说明}：本节只给最终答案便于快速核对，完整推导见解析版。答案为编写者独立推导，\textbf{非官方标准答案}。
\end{gknotice}

\gkpart{选择题}
\begin{center}\begin{tabular}{ccccccccccc}
\toprule
题号 & 1 & 2 & 3 & 4 & 5 & 6 & 7 & 8 & 9 & 10 \\
\midrule
答案 & B & A & C & A & C & D & B & C & D & D \\
\bottomrule
\end{tabular}\end{center}

\gkpart{填空题}
\begin{itemize}[leftmargin=2.6em]
  \item[\textbf{11.}] 存在 $x\in\mathbf R$，使 $\lvert x-2\rvert+\lvert x-4\rvert\leqslant3$。
  \item[\textbf{12.}] $15$。
  \item[\textbf{13.}] $4$。
  \item[\textbf{14.}] $10$。
  \item[\textbf{15.}] ②④。
\end{itemize}

\gkpart{解答题}
\begin{itemize}[leftmargin=2.6em]
  \item[\textbf{16.}] （1）$A=\dfrac\pi3$；（2）$b=4$，$c=6$。
  \item[\textbf{17.}] （1）减区间 $(-\infty,\ln2)$，增区间 $(\ln2,+\infty)$，极小值 $2a+2-2\ln2$，无极大值；（2）证明见解析。
  \item[\textbf{18.}] （1）证明；（2）证明；（3）$60^\circ$。
  \item[\textbf{19.}] （1）$\dfrac{x^2}{16}+\dfrac{y^2}{12}=1$；（2）$2x-y-1=0$；（3）不存在关于 $l$ 对称的相异两点。
  \item[\textbf{20.}] 证明见解析。
  \item[\textbf{21.}] （1）$\{0,2,4,6,8\}$；（2）$P(X=0)=\dfrac1{24}$，$P(X=2)=\dfrac18$，$P(X=4)=\dfrac7{24}$，$P(X=6)=\dfrac38$，$P(X=8)=\dfrac16$；（3）① $\dfrac1{216}$，② 该品酒师有较强的酒味鉴别能力。
\end{itemize}
